% Generated draft: requires mathematical and source review.
\documentclass[11pt]{article}
\usepackage[margin=25mm]{geometry}
\usepackage{fontspec}
\setmainfont{Latin Modern Roman}
\usepackage{amsmath,amssymb,mathtools,unicode-math}
\setmathfont{Latin Modern Math}
\usepackage{enumitem,needspace,xurl,hyperref}
\setlength{\parindent}{0pt}
\setlength{\parskip}{5pt}
\setlength{\emergencystretch}{3em}
\newcommand{\N}{\mathbb{N}}
\newcommand{\Z}{\mathbb{Z}}
\newcommand{\Q}{\mathbb{Q}}
\newcommand{\R}{\mathbb{R}}
\newcommand{\C}{\mathbb{C}}
\newcommand{\F}{\mathbb{F}}
\newcommand{\A}{\mathbb{A}}
\newcommand{\PP}{\mathbb{P}}
\newcommand{\E}{\mathbb{E}}
\newcommand{\eps}{\varepsilon}
\newcommand{\dd}{\,\mathrm{d}}
\begin{document}
\section*{UnsolvedMath ID: 30001445}
\textbf{Source problem code:} catalogue entry
\subsection{Definitions and mathematical statement}
The supplied record poses a question about finite quotients of a finitely generated residually finite group, but does not specify the quotient family or its coarse metric. Its background additionally describes the extracted passage as an incomplete continuation concerning asymptotic assembly maps. Thus the following gives a conventional interpretation of the displayed question, conditional on the qualifications below; it does not identify the intended original formulation conclusively.

A group \(\Gamma\) is residually finite if every nonidentity element has nonidentity image in some finite quotient. Choose a finite symmetric generating set \(S\). Under the group-quotient interpretation, let \((\Gamma_i)_{i\in\N}\) be a family of finite-index normal subgroups of \(\Gamma\), and give \(\Gamma/\Gamma_i\) the quotient word metric induced by \(S\). Write
\[
X(\Gamma)=\bigsqcup_{i\in\N}\Gamma/\Gamma_i.
\]
For the usual coarse interpretation, equip this union with a proper metric restricting to these quotient metrics and separating its components: for every \(R>0\), only finitely many pairs of distinct components have distance at most \(R\).

Property \(\tau\) relative to this family means a uniform spectral gap in the quotient representations. Explicitly, if \(\lambda_i\) is the left translation representation on \(\ell^2(\Gamma/\Gamma_i)\), there is \(\varepsilon>0\) such that
\[
\sum_{s\in S}\|\lambda_i(s)\xi-\xi\|_2^2
\geq \varepsilon\|\xi\|_2^2
\quad\text{for every }i\in\N
\text{ and every }\xi\perp\mathbf{1}.
\]
The record does not say whether its unqualified property \(\tau\) means this relative property or the property with respect to all finite-index normal subgroups.

For a discrete metric space \(X\) as above, let \(P_R(X)\) be the Rips complex whose simplices are the finite subsets of diameter at most \(R\), and let \(K_j^{\mathrm{lf}}\) denote locally finite analytic \(K\)-homology. In the usual maximal coarse Baum--Connes formulation, put
\[
KX_j(X)=\varinjlim_{R\to\infty}K_j^{\mathrm{lf}}(P_R(X)),
\qquad j\in\{0,1\}.
\]
Let \(C^*_{\max}(X)\) denote the maximal Roe algebra, namely the maximal \(C^*\)-completion of the algebra of locally compact finite-propagation operators on an ample \(X\)-module. Here an ample \(X\)-module is a Hilbert space with a nondegenerate representation of \(C_0(X)\) in which no nonzero function acts compactly; local compactness means that multiplication on either side by any function in \(C_0(X)\) gives a compact operator.

Under these conventions, the displayed question asks whether, for every such \(\Gamma\) and quotient family satisfying the intended property \(\tau\) hypothesis, the maximal coarse assembly homomorphisms
\[
\mu^{\max}_{X(\Gamma),j}\colon
KX_j(X(\Gamma))\longrightarrow
K_j\!\left(C^*_{\max}(X(\Gamma))\right)
\]
are isomorphisms for both \(j=0\) and \(j=1\).

The source does not specify whether the subgroups must be nested, have trivial intersection, or have indices tending to infinity, nor whether normality is intended instead of a Schreier-quotient interpretation. No such additional condition is asserted here as an established source hypothesis. The background's reference to asymptotic assembly-map properties may indicate a different intended target from the ordinary maximal coarse assembly maps displayed above. These omissions require source review before this conditional formulation can be accepted as a precise transcription.
\subsection{Short English statement}
The supplied question asks whether the maximal coarse Baum--Connes assembly map is an isomorphism in both degrees for the coarse disjoint union of finite quotients of a finitely generated residually finite group with property \(\tau\). The quotient-family hypotheses, the meaning of property \(\tau\), and the possible intended asymptotic version remain unresolved.
\subsection{Sources}
\begin{itemize}
\item Local UnsolvedMath registry, permanent ID \(30001445\), source problem code \texttt{catalogue entry}, title ``Maximal Coarse Baum--Connes Assembly for Expander Quotients.'' The supplied record belongs to the Oberwolfach Reports Open Problems collection and gives the proposed year as 2010; no individual proposer is identified.
\item Supplied report citation: \emph{C*-Algebren} (2010), Oberwolfach Reports, pp.\ 667--732, DOI \texttt{10.4171/owr/2010/13}. Source URL and literature-source link labelled ``original OWR report'': \url{https://doi.org/10.4171/owr/2010/13}. No exact passage or problem page is supplied.
\item Supplied literature-source link labelled ``Oyono-Oyono--Yu, maximal Roe algebras of expanders'': \url{https://doi.org/10.1016/j.jfa.2009.04.017}.
\item Literature assessment dated \texttt{2026-08-21}: the supplied curation labels the problem open, describes the extracted passage as an incomplete continuation concerning asymptotic assembly-map properties for expander families, and states that existing work proves the maximal coarse Baum--Connes map for important classes rather than the general target suggested here. This is supplied dated curation, not an independently verified literature assessment; its status and cited claims require independent verification.
\item The registry's supplied statement-status label is \texttt{corrected\_verified}. Its verification note says only that TeX notation was enclosed in math delimiters without changing mathematical content; this does not resolve the formulation issues above.
\end{itemize}
\end{document}
